% One-time imported draft. Editable native source; no Markdown regeneration.
\section{Step 20 Render and publish the newly current
  board}\label{008-life-generations-s20-step-20-render-and-publish-the-newly-current-board}

The back-edge returns to the renderer. After rendering, the same ready
label is reached with the new generation. The boundary test then decides
whether another sweep is required.

\subsection{Write or
  inspect}\label{008-life-generations-s20-write-or-inspect}

Append the loop back-edge.

\MARSCodeLabel{MIPS assembly}
\begin{minted}{mars-idea}
    beq $zero, $zero, render
    nop
\end{minted}

\subsection{Expected
  observation}\label{008-life-generations-s20-expected-observation}

The complete default program publishes generation zero, one, and two
before exiting.

\subsection{Check your
  understanding}\label{008-life-generations-s20-check-your-understanding}

Why reuse the renderer after swapping?

\begin{itemize}
  \tightlist
  \item
        \textbf{A} It always reads the current pointer, whichever physical
        board now owns that role.
  \item
        \textbf{B} The frame automatically follows pointer assignments.
  \item
        \textbf{C} It recalculates neighbors as part of every pixel store.
\end{itemize}
